从\hyperref[eq:5-6]{式(5-6)}可立即得到假设 $H_0: \boldsymbol{\mu}=\boldsymbol{\mu}_0$, $H_1: \boldsymbol{\mu}\neq\boldsymbol{\mu}_0$ 的一个检验方法. 若在显著性水平 $\alpha$ 下,

\phantomsection\label{eq:5-7}
\begin{equation}
T^2=n(\bar{\mathbf{x}}-\boldsymbol{\mu}_0)'\mathbf{S}^{-1}(\bar{\mathbf{x}}-\boldsymbol{\mu}_0)>\frac{(n-1)p}{(n-p)}F_{p,n-p}(\alpha) \tag{5-7}
\end{equation}

则拒绝 $H_0$ 而接受 $H_1$.

简要讨论一下 $T^2$ 分布的实质及其与一元检验统计量之间的对应关系, 有助于我们对 $T^2$ 分布的了解. 在 4.4 节中我们曾指出, 威沙特分布是 $\chi^2$ 分布的推广, 由于

\[
T^2=\sqrt{n}(\bar{\mathbf{X}}-\boldsymbol{\mu}_0)'\left(\frac{\sum_{j=1}^{n}(\mathbf{X}_j-\bar{\mathbf{X}})(\mathbf{X}_j-\bar{\mathbf{X}})'}{n-1}\right)^{-1}\sqrt{n}(\bar{\mathbf{X}}-\boldsymbol{\mu}_0)
\]

通过关系

\phantomsection\label{eq:5-8}
\begin{equation}
\begin{aligned}
T^2_{p,n-1}&=\begin{pmatrix}\text{多元正态}\\ \text{随机向量}\end{pmatrix}'\begin{pmatrix}\text{威沙特随机矩阵}\\ \text{自由度}\end{pmatrix}^{-1}\begin{pmatrix}\text{多元正态}\\ \text{随机向量}\end{pmatrix}\\
&=N_p(\mathbf{0},\boldsymbol{\Sigma})'\left[\frac{1}{n-1}W_{p,n-1}(\boldsymbol{\Sigma})\right]^{-1}N_p(\mathbf{0},\boldsymbol{\Sigma})
\end{aligned} \tag{5-8}
\end{equation}

与正态随机向量 $N_p(\mathbf{0},\boldsymbol{\Sigma})$ 及威沙特随机矩阵 $W_{p,n-1}(\boldsymbol{\Sigma})$ 相联系, 所以, $T^2$ 与一元情形下的

\[
t^2=\sqrt{n}(\bar{X}-\mu_0)(s^2)^{-1}\sqrt{n}(\bar{X}-\mu_0)
\]

或

\[
t^2_{n-1}=\left(\text{正态随机变量}\right)\left(\frac{\chi^2\text{ 随机变量}}{\text{自由度}}\right)^{-1}\left(\text{正态随机变量}\right)
\]

类似. 由于多元正态随机变量的分布与威沙特随机变量的分布相互独立[见式(4-23)], 因此它们的联合密度函数等于边缘正态分布与威沙特分布之积. 利用这个联合分布及表达\hyperref[eq:5-8]{式(5-8)}, 通过积分可推导出\hyperref[eq:5-5]{式(5-5)} $T^2$ 的分布.

在多元情形中, 人们很少满足于检验假设 $H_0: \boldsymbol{\mu}=\boldsymbol{\mu}_0$, 在这个零假设中, 均值向量的所有分量都是完全确定的. 人们通常更倾向于按照观测数据寻求一个 $\boldsymbol{\mu}$ 的似真值区域. 我们将在 5.4 节中再讨论这个问题.

\noindent\textbf{例 5.1}\quad ($T^2$ 的计算)

设一个容量为 $n=3$ 的随机样本取自二维正态总体, 其数据矩阵为

\[
\mathbf{X}=\begin{bmatrix}6&9\\10&6\\8&3\end{bmatrix}
\]

试对 $\boldsymbol{\mu}'=[9,5]$ 计算 $T^2$ 的观测值. 此时 $T^2$ 的抽样分布是什么? 注意到

\[
\bar{\mathbf{x}}=\begin{bmatrix}\bar{x}_1\\ \bar{x}_2\end{bmatrix}=\begin{bmatrix}\frac{6+10+8}{3}\\ \frac{9+6+3}{3}\end{bmatrix}=\begin{bmatrix}8\\6\end{bmatrix}
\]

且

\begin{align*}
s_{11}&=\frac{(6-8)^2+(10-8)^2+(8-8)^2}{2}=4\\
s_{12}&=\frac{(6-8)(9-6)+(10-8)(6-6)+(8-8)(3-6)}{2}=-3\\
s_{22}&=\frac{(9-6)^2+(6-6)^2+(3-6)^2}{2}=9
\end{align*}
